\begin{enumerate}
	\item
	\begin{enumerate}
		\item Un nombre premier possède exactement deux diviseurs : 1 et lui-même.

		1 n'est pas premier, $4 = 2\times 2$ et $6 = 2\times 3$ ne sont pas premiers.

		Les boules portant un numéro premier sont les boules 2, 3, 5 et 7.

		\item Les boules sont indiscernables au toucher donc le tirage est équiprobable. Il y a 4 boules favorables sur 7 boules au total.
		\[p = \dfrac{4}{7}\]
		
	La probabilité d'obtenir une boule dont le numéro est un nombre premier est de $\dfrac{4}{7}$.
	\end{enumerate}

	\item
	\begin{enumerate}
		\item Par lecture graphique, la fréquence d'apparition de la boule numéro 2 est de 10\,\%.

		\item $10 + 12 + 22 + 10 = 54$.

		La somme des fréquences des numéros premiers est bien égale à 54\,\%.
	\end{enumerate}

	\item
	\begin{enumerate}
		\item $\dfrac{4}{7}\approx 0,571$ alors que $54\,\% = \dfrac{54}{100}=0,54$.

		Donc $\dfrac{4}{7}$ n'est pas égal à 54\,\%.

		\item La probabilité est une valeur théorique, calculée à partir de la composition de la boîte ; la fréquence est une valeur observée sur un nombre fini de tirages (100 ici) et varie d'une série d'expériences à l'autre.
		La fréquence se rapproche de la probabilité lorsque le nombre de tirages devient très grand. Or répéter 100 fois l'expérience n'est pas suffisant pour que la fréquence observée donne la même valeur que la probabilité.
	\end{enumerate}
\end{enumerate}
